%% ============================================================
%%  走れメロス ── 日本風表紙「藍と朱の疾走」
%%  エンジン: LuaLaTeX (TeX Live 2026, Windows)
%%  standalone クラス使用（ページを tikzpicture に合わせる）
%%
%%  フォント優先順位:
%%    1. 游明朝 (Yu Mincho)   ← Windows 8.1 以降に標準搭載
%%    2. MS 明朝              ← 旧 Windows でも利用可
%%  どちらも無い場合は TeX Live 付属の IPAex 明朝を使用
%% ============================================================
\documentclass[border=0pt]{standalone}
\usepackage{luatexja-fontspec}   %% LuaLaTeX 用 日本語フォント指定
\usepackage{tikz}
\usetikzlibrary{calc}
\usepackage{xcolor}

%% 游明朝が入っていれば游明朝、なければ MS 明朝にフォールバック
%% （どちらも Windows 標準搭載）
\setmainjfont[
  BoldFont = Yu Mincho Demibold,   %% 游明朝 Demibold
]{Yu Mincho}                        %% 游明朝 Regular

%% ── 游明朝が見つからない場合のフォールバック例 ──────────────
%%   上の \setmainjfont 行を以下に差し替えてください
%%   \setmainjfont{MS Mincho}       %% MS 明朝
%%   \setmainjfont{IPAexMincho}     %% TeX Live 付属（要 ipaex パッケージ）

\definecolor{aiiro}   {HTML}{1A2B55}
\definecolor{aiirodk} {HTML}{0F1A38}
\definecolor{shuiro}  {HTML}{B83030}
\definecolor{kiniro}  {HTML}{C9A347}
\definecolor{washi}   {HTML}{F2E4C6}
\definecolor{washidk} {HTML}{C8B890}

%% ページ中心 (45,80)mm、大円中心 (45,88)mm
\newcommand\CX{45}
\newcommand\CY{80}
\newcommand\CCX{45}
\newcommand\CCY{88}

\begin{document}
\begin{tikzpicture}[x=1mm, y=1mm]

  \clip (0,0) rectangle (90,160);

  %% ---- 1. 背景 -----------------------------------------------
  \fill[aiiro] (0,0) rectangle (90,160);
  \fill[aiirodk, opacity=0.55] (0,0) rectangle (90,38);
  \fill[aiirodk, opacity=0.40] (0,128) rectangle (90,160);

  %% ---- 2. 青海波（せいがいは） r=5mm -----------------------
  \begin{scope}[opacity=0.09, draw=washi, fill=none, line width=0.45pt]
    \foreach \row in {-18,...,18} {
      \pgfmathsetmacro{\xoff}{0.5*(1 - cos(\row * 180)) * 5}
      \foreach \col in {-11,...,11} {
        \draw ({(\col*10+\xoff)+\CX},{(\row*5)+\CY})
              arc (180:0:5);
      }
    }
  \end{scope}

  %% ---- 3. 朱色の大円 ----------------------------------------
  \fill[shuiro, opacity=0.16] (\CCX,\CCY) circle (42);
  \fill[shuiro, opacity=0.10] (\CCX,\CCY) circle (46);
  \fill[shuiro]               (\CCX,\CCY) circle (30);
  \fill[shuiro!65!white, opacity=0.22] (\CCX,{(\CCY)+8}) circle (16);
  \draw[kiniro, line width=1.3pt]  (\CCX,\CCY) circle (31.5);
  \draw[kiniro, line width=0.45pt] (\CCX,\CCY) circle (33.2);
  \draw[kiniro, line width=0.18pt] (\CCX,\CCY) circle (35);

  %% ---- 4. タイトル「走れメロス」縦積み ----------------------
  \foreach \ch/\i in {走/0, れ/1, メ/2, ロ/3, ス/4} {
    \pgfmathsetmacro{\yy}{(\CCY)+(22-\i*11)}
    \node[text=washi] at (\CCX,\yy)
      {\fontsize{26pt}{26pt}\selectfont\ch};
  }

  %% ---- 5. 金色の水平帯 --------------------------------------
  \fill[kiniro, opacity=0.13] (8,43) rectangle (82,49);
  \draw[kiniro, line width=0.7pt]  (8,49) -- (82,49);
  \draw[kiniro, line width=0.22pt] (8,43) -- (82,43);

  %% ---- 6. 作者名「太宰　治」縦積み --------------------------
  \foreach \ch/\i in {太/0, 宰/1, 治/2} {
    \pgfmathsetmacro{\yy}{41-\i*11}
    \node[text=washidk] at (\CX,\yy)
      {\fontsize{15pt}{15pt}\selectfont\ch};
  }

  %% ---- 7. 家紋風モチーフ（上部中央）「丸に十字」 -----------
  \draw[kiniro, line width=0.85pt] (\CX,144) circle (5);
  \draw[kiniro, line width=0.35pt] (\CX,144) circle (3.6);
  \draw[kiniro, line width=0.5pt]  ({(\CX)-5},144) -- ({(\CX)+5},144);
  \draw[kiniro, line width=0.5pt]  (\CX,139) -- (\CX,149);
  \foreach \dx in {9,13,17,21} {
    \draw[kiniro, line width=0.22pt]
      ({(\CX)+\dx},144) -- ({(\CX)+\dx+4.5},144);
    \draw[kiniro, line width=0.22pt]
      ({(\CX)-\dx},144) -- ({(\CX)-\dx-4.5},144);
  }

  %% ---- 8. 菱形装飾（下部中央）--------------------------------
  \fill[kiniro, opacity=0.85]
    (\CX,13.8) -- ({(\CX)+3.8},10) -- (\CX,6.2) -- ({(\CX)-3.8},10) -- cycle;
  \draw[kiniro!70!white, line width=0.3pt]
    (\CX,15.5) -- ({(\CX)+5.5},10) -- (\CX,4.5) -- ({(\CX)-5.5},10) -- cycle;
  \foreach \dx in {9,13,17,21} {
    \draw[kiniro, line width=0.22pt]
      ({(\CX)+\dx},10) -- ({(\CX)+\dx+4.5},10);
    \draw[kiniro, line width=0.22pt]
      ({(\CX)-\dx},10) -- ({(\CX)-\dx-4.5},10);
  }

  %% ---- 9. 二重枠線（金色）-----------------------------------
  \draw[kiniro, line width=1.4pt] (5,5) rectangle (85,155);
  \draw[kiniro, line width=0.3pt] (7,7) rectangle (83,153);

  %% ---- 10. 四隅の角飾り--------------------------------------
  \draw[kiniro, line width=0.6pt] (5,11) -- (5,5) -- (11,5);
  \draw[kiniro, line width=0.6pt] (79,5) -- (85,5) -- (85,11);
  \draw[kiniro, line width=0.6pt] (5,149) -- (5,155) -- (11,155);
  \draw[kiniro, line width=0.6pt] (79,155) -- (85,155) -- (85,149);

\end{tikzpicture}
\end{document}